% ============================================================================
\chapter{Pravci daljnjeg istra\v{z}ivanja}
\label{ch:pravci}
% ============================================================================

U ovom poglavlju izdvajaju se mogu\'ci pravci daljnjeg istra\v{z}ivanja koji
prirodno proizlaze iz ura\dj{}enog rada. Svaki pravac povezan je sa konkretnom
prazninom uo\v{c}enom u radu ili u pregledu literature iz poglavlja
\ref{ch:srodni}. Pravci su pore\dj{}ani od onih koji najvi\v{s}e doprinose
nau\v{c}noj vrijednosti rada prema onima koji predstavljaju \v{s}ire pro\v{s}irenje.
Cilj poglavlja je pokazati da postoje\'ci rezultati \v{c}ine osnovu za vi\v{s}e
zaokru\v{z}enih istra\v{z}iva\v{c}kih cjelina.

\section{Kvantifikacija tro\v{s}ka operativnog sistema}
\label{sec:pravci-overhead}

Najva\v{z}niji pravec jeste mjerenje tro\v{s}ka koji unosi operativni sistem u
odnosu na golo izvr\v{s}avanje. Ista operacija bi se izmjerila sa raspore\dj{}iva\v{c}em
i bez njega, \v{c}ime bi se dobila stvarna cijena upotrebe operativnog sistema na
ovoj mekoj jezgri. Ovaj pravec direktno odgovara na pitanje cijene koje je
\v{c}esto u sredi\v{s}tu literature o operativnim sistemima za rad u realnom
vremenu \citep{ulsolutions2022rtos}.

\section{Uskla\dj{}ivanje metrike i statisti\v{c}ka obrada}
\label{sec:pravci-metrika}

Mjerenje iz poglavlja \ref{ch:mjerenje} obuhvata i tro\v{s}ak rada sa semaforima.
Da bi se rezultat mogao pore\dj{}ivati sa literaturom, potrebno je dodatno
izmjeriti u\v{z}u, \v{c}istu zamjenu konteksta bez tog tro\v{s}ka. Uz to, umjesto
najmanje, prosje\v{c}ne i najve\'ce vrijednosti, rezultate bi trebalo izraziti
srednjom vrijedno\v{s}\'cu sa standardnom devijacijom i raspodjelom, te
karakterizovati podrhtavanje. Time bi mjerenja dobila statisti\v{c}ku strogost
koju nau\v{c}ni rad zahtijeva.

\section{Uticaj mikroarhitekture i skupa instrukcija}
\label{sec:pravci-mikroarhitektura}

Meka jezgra dopu\v{s}ta promjenu konfiguracije, na primjer uklju\v{c}ivanje
sa\v{z}etog skupa instrukcija ili sklopa za pomjeranje. Mjerenjem zamjene
konteksta pri razli\v{c}itim konfiguracijama dobili bi se trendovi koji povezuju
mikroarhitekturu sa cijenom operativnog sistema. Ovaj pravec je u skladu sa
istra\v{z}ivanjima koja zajedno posmatraju mikroarhitekturu i operativni sistem
\citep{scheck2026coexploration}.

\section{Latencija i determinizam prekida pod optere\'cenjem}
\label{sec:pravci-prekidi}

Pored zamjene konteksta, va\v{z}na veli\v{c}ina je latencija prekida i njena
postojanost pod razli\v{c}itim optere\'cenjem sistema. Mjerenjem latencije
prekida i podrhtavanja, uz vi\v{s}e zadataka i ve\'ce optere\'cenje, ocijenila bi
se predvidljivost sistema, \v{s}to je klju\v{c}no svojstvo sistema za rad u
realnom vremenu.

\section{Sistematsko pore\dj{}enje simulatora i hardvera}
\label{sec:pravci-qemu}

Po\v{s}to je emulator QEMU funkcijski ta\v{c}an, a ne takovno ta\v{c}an
\citep{qemu}, zanimljivo je sistematski ispitati koliko dobro on predvi\dj{}a
vremenske veli\v{c}ine na stvarnom hardveru. Pore\dj{}enjem vi\v{s}e mjera na obje
platforme dobio bi se sud o tome u kojim slu\v{c}ajevima je simulacija dovoljno
vjerodostojna, a u kojima nije.

\section{Hardversko ubrzanje raspore\dj{}ivanja}
\label{sec:pravci-hardver}

Kona\v{c}no, budu\'ci da je rije\v{c} o mekoj jezgri na programabilnoj logici,
otvara se mogu\'cnost dodavanja namjenskog sklopa za ubrzanje raspore\dj{}ivanja
ili zamjene konteksta. Takav pristup je u literaturi pokazao zna\v{c}ajna
ubrzanja \citep{balas2023cv32rt}, a ovaj rad pru\v{z}a polaznu osnovu i mjerni
okvir naspram kojeg bi se takvo ubrzanje moglo vrednovati.

\section{Zaklju\v{c}ak poglavlja}
\label{sec:pravci-zakljucak}

U poglavlju je izlo\v{z}eno \v{s}est pravaca daljnjeg istra\v{z}ivanja, od
kvantifikacije tro\v{s}ka operativnog sistema i statisti\v{c}ke obrade mjerenja,
preko uticaja mikroarhitekture i pona\v{s}anja prekida, do sistematskog
pore\dj{}enja sa simulatorom i hardverskog ubrzanja raspore\dj{}ivanja. Svaki
pravec se naslanja na postoje\'ce rezultate i prazninu u literaturi, \v{c}ime se
pokazuje da rad ima jasan potencijal za nastavak.
